% Standalone section: requires amsmath, amssymb, and hyperref (for \url).
% No custom macros or theorem environments are required.
\section{Categorical consequences of the rank-$26$ reconstruction}
\label{sec:r155-new-categorical}

The results below give a construction obstruction and restrictions on
algebras in the newly reconstructed fusion ring. They do not establish or
exclude a categorification of $R_6^{4,155}$. For the rank-$26$ category
$\mathcal B$, we use the numbering of the accompanying reconstruction:
the unit is object $1$, the fermion is object $0$, and
\[
 d=4+3\sqrt2,\qquad b=2+4d=18+12\sqrt2,\qquad
 A=23,\quad L=25,
\]
where $\operatorname{FPdim}(A)=b-1$, $\operatorname{FPdim}(L)=b$, and
$\operatorname{FPdim}(\mathcal B)=8b^2$. All assertions about
$\mathcal B$ are conditional on a positive ribbon categorification of
the reconstructed data.

\subsection{The new order-$32$ near-group centers cannot give the target
by dimension-eight condensation}

Radchenko--Wheeler construct complex near-group categories of type
$(G,|G|)$ for every product of two cyclic groups; unitarity is not
asserted. Their primary source is
\emph{Real quadratic fields and finite quantum dilogarithms I},
\url{https://arxiv.org/abs/2609.21892}, version~2, Section~5.4,
Theorem~12 and Corollary~6. Thus their construction supplies the three
order-$32$ groups $C_{32}$, $C_2\times C_{16}$, and $C_4\times C_8$.
The following argument applies directly to their algebraic construction;
it does not import a classification of unitary centers.

\paragraph{A half-braiding lemma.}
In the strict endomorphism realization, write the simple objects as
$\alpha_g$ and $\rho$. The displayed identities in the cited Section~5.4
include
\[
 \alpha_g\rho=\rho,\quad
 \rho\alpha_g=\operatorname{Ad}(U_g)\rho,\quad
 U_gU_h=U_{g+h},\quad
 \alpha_g(U_h)=\chi(g,h)U_h,
\]
where $\chi$ is a nondegenerate symmetric bicharacter. There is also a
nonzero split inclusion $T:\rho\longrightarrow\rho^2$.
If a central object $Z$ has multiplicity-free pointed underlying object
$\bigoplus_{g\in H}\alpha_g$, then every one of its summands is central,
with self-braiding $\chi(g,g)^{-1}$.

\emph{Proof.}
Its half-braiding with $\alpha_h$ is diagonal, with scalar $a(g,h)$ on
label $g$. Write the entry of its half-braiding with $\rho$ from label
$g$ to label $k$ as $M_{kg}U_k$. Naturality at
$U_h:\rho\to\rho\alpha_h$ and the hexagon give
\[
 U_hM_{kg}U_k=a(k,h)M_{kg}U_k\alpha_g(U_h).
\]
Thus $M_{kg}\ne0$ implies $a(k,h)\chi(g,h)=1$ for every $h$.
Nondegeneracy of $\chi$ permits at most one nonzero entry in each row
of $M$. Since $M$ is invertible, it is monomial, with a permutation
$\pi$ of the labels $H$. The hexagon for $\rho^2$ has output label
$\pi^2(g)$, whereas naturality at $T$ compares it with a map having
output label $\pi(g)$. Both maps are nonzero, being compositions of
isomorphisms and split inclusions. Projection onto the tagged direct
summands therefore gives $\pi^2(g)=\pi(g)$. Hence $\pi$ is the identity.
All half-braidings preserve each summand, and the preceding scalar
identity gives $a(g,h)=\chi(g,h)^{-1}$.\hfill$\square$

\paragraph{Proposition.}
None of these three order-$32$ Radchenko--Wheeler categories has a
connected \'etale algebra of dimension eight in its center. Consequently
their centers cannot produce the proposed center of $R_6^{4,155}$ by a
single connected-\'etale condensation.

\emph{Proof.}
Use the established two-dimension reduction to pass to a positive
spherical Galois representative; a hypothetical algebra transports under
the same conjugation. Put
\[
 r=16+12\sqrt2=4d,\qquad
 \Delta=32+r^2=8(72+48\sqrt2).
\]
The formal-codegree decomposition of $I(1)$ has constituent dimensions
$1$, $r+1$, and $31$ copies of $r+2$. Indeed the $31$ nontrivial group
characters have formal codegree $32$, while the two remaining characters
send $\rho$ to its two quadratic roots. In particular every nonunit
constituent of $I(1)$ has dimension greater than eight.

Suppose $E\in Z(\mathcal C)$ is connected \'etale of dimension eight.
Its underlying object is pointed, since $r>8$. Adjunction and the
preceding constituent bound give
\[
 \dim\operatorname{Hom}_{\mathcal C}(1,F(E))
 =\dim\operatorname{Hom}_{Z(\mathcal C)}(I(1),E)=1.
\]
Thus $F(E)$ is a connected separable algebra in the pointed subcategory.
It is simple as a right module over itself, since its endomorphism space
is $\operatorname{Hom}_{\mathcal C}(1,F(E))$. Each $\alpha_g\otimes F(E)$
is also simple. Free-module adjunction now shows
\[
 [F(E):\alpha_g]
 =\dim\operatorname{Hom}_{F(E)}(\alpha_g\otimes F(E),F(E))\in\{0,1\}.
\]
Its support is the stabilizer subgroup $H$ of this simple module and has
order eight. The lemma splits $E$ into central invertibles. Twist-freeness
of a connected \'etale algebra then forces $\chi(h,h)=1$ for every
$h\in H$.

For completeness the required bicharacter count is elementary and
uniform over all nondegenerate choices. On $C_{32}$,
$\chi(x,x)=\exp(2\pi i c x^2/32)$ with $c$ odd, so the bosons are exactly
the four multiples of eight. On $C_2\times C_{16}$, write diagonal
parameters $p/2,q/16$ and off-diagonal parameter $r_0/2$.
Nondegeneracy forces $p=1$ and $q$ odd: otherwise the order-$16$
generator maps to a character of smaller order, or the order-two
generator is in the pairing kernel or has the same character as eight
times that generator. The boson equation is
$8a+qb_0^2=0\pmod {16}$; it forces $a=0$ and $4\mid b_0$.
On $C_4\times C_8$, write parameters $p/4,q/8,r_0/4$.
Here $p,q$ must be odd: if $q$ is even the order-eight character has
smaller order, while if $p$ is even the nonzero element $(2,4r_0)$ is in
the kernel. The boson equation
\[
 2pa^2+qb_0^2+4r_0ab_0=0\pmod8
\]
first forces $b_0$ even, then $a$ even, and finally $4\mid b_0$.
There are again only four bosons. No required subgroup $H$ exists.
Finally, the source-center dimension is $\Delta^2=64\dim(\mathcal C_{155})^2$;
the local-module dimension formula forces a condensation to the target
center to use an algebra of dimension eight.\hfill$\square$

This excludes all dimension-eight connected \'etale algebras in these
centers, not merely pointed candidates. It does not exclude other source
categories, products followed by condensation, or the unresolved
near-group type $(8,16)$.

\subsection{The only possible connected \'etale underlying objects in
$\mathcal B$}

\paragraph{Proposition.}
Every connected \'etale algebra in a positive ribbon categorification of
the reconstructed rank-$26$ data has underlying object
\[
             1,\qquad 1+A,\qquad 1+A+L.
\]
This statement does not assert uniqueness of algebra structures.

\emph{Proof.}
The twist-one simples are precisely $1,2,5,23,25$, with $2^*=5$.
The reconstructed fusion rules give
\[
 2\otimes2=9+10+16+19,\qquad
 2\otimes5=1+8+11+22.
\]
The first product has no twist-one constituent; the second has only the
unit among such constituents. A connected \'etale algebra has a
nondegenerate Frobenius pairing. If it contained object $2$, it would
contain $5$, and multiplication of suitable copies of $2$ and $5$ to the
unit would be nonzero. Its multiplication on $2\otimes2$ would vanish.
Associativity on $2\otimes2\otimes5$ is then impossible: multiplying the
first two factors gives zero, whereas contracting the last two to the
unit and multiplying by the first gives a nonzero map. This argument
also works with arbitrary multiplicities, by choosing a nonzero entry
of the pairing on the multiplicity spaces. Thus neither $2$ nor $5$
can occur.

Write the remaining candidate as $E=1+sA+tL$, with $s,t\geq0$.
Its dimension is $1-s+(s+t)b$, at most
$\sqrt{\operatorname{FPdim}(\mathcal B)}=2\sqrt2\,b$.
If $s+t\geq3$, its dimension is at least $3b-2$, which exceeds this
bound by four. Hence $s+t\leq2$. Apart from the three stated objects,
the possibilities have dimensions $b+1$, $2b-1$, and $2b+1$, with norms
\[
 N(b+1)=73,\qquad N(2b-1)=73,\qquad N(2b+1)=217.
\]
But
\[
 N(\operatorname{FPdim}(\mathcal B))=N(8b^2)
   =82944=2^{10}3^4.
\]
The number $\operatorname{FPdim}(\mathcal B)/\operatorname{FPdim}(E)^2$
must be an algebraic integer, being the dimension of the local-module
category. Its norm would require $73^2$ or $217^2$ to divide $82944$,
which is false.\hfill$\square$

The standard algebra facts used here---the Frobenius pairing,
twist-freeness in the positive setting, and local-module dimension---are
discussed in Davydov--M\"uger--Nikshych--Ostrik,
\emph{The Witt group of non-degenerate braided fusion categories},
Section~3, \url{https://arxiv.org/abs/1009.2117}.

\paragraph{Fusion subcategories and restriction.}
There are exactly four fusion subcategories: $\mathrm{Vec}$,
$\mathrm{sVec}$, the rank-$16$ fermion centralizer, and $\mathcal B$.
This conclusion follows directly from a few support properties of the
full fusion table. Every noninvertible simple other than $4,7$ has $22$
in its product with its dual. For $4,7$, that product contains $8$, whose
product with its dual contains $22$. The square of $22$ contains every
even simple except the fermion; it contains $23$, and $22\otimes23$
contains the fermion as well. Thus every noninvertible even simple
generates the entire even category. A subcategory containing it and an
odd object is faithfully $C_2$-graded with that even component. Equality
of the dimensions of grading components forces it to be all of
$\mathcal B$.
The accompanying integer checker also verifies every one of the $676$
restriction-product identities
\[
 W(i)W(j)=\sum_k N_{ij}^{k}W(k)
\]
in the original rank-six ring. Thus the full reconstructed fusion rules
pass the restriction-functor consistency check, not only its previously
available low products. The checks and exact support sets are recorded
in \texttt{full\_fusion\_categorical\_receipt.json}.

\subsection{The canonical intermediate condensation adds no new
near-group category}

When $\mathcal B$ comes from a categorification $\mathcal C$ of the
rank-six ring, the restricted adjoint algebra is $E=1+A+L$.
Its even subalgebra $E_0=1+A$ has dimension $b$.
The local category $\mathcal K=\mathcal B_{E_0}^{0}$ has dimension eight.
The summand $L$, as an $E_0$-module, is an invertible boson $\ell$;
the free fermion $f_0=fE_0$ is distinct, and $f_0\ell\ne\ell$ has the
same underlying object $L$. Their twists $1,-1,1,1$ give a toric-code
factor. Preservation of central charge then gives
\[
 \mathcal K\simeq\mathrm{ToricCode}\boxtimes\mathrm{Semion}^{\mathrm{rev}}.
\]

In fact the category of all modules $\mathcal D=\mathcal B_{E_0}$ is
tensor equivalent to $\mathcal C\boxtimes\operatorname{Rep}(C_2)$.
Here are details preventing an inference from merely fixing simple
objects. The boson $\ell$ realizes $\mathcal D$ as a
$C_2$-equivariantization of $\mathcal C$. Its action fixes $X,Y$, since
each is the free image of a $d$-dimensional simple of $\mathcal B$.
It fixes every invertible as well, since their unique central lifts
have distinct twists $1,-1,i,-i$. Consequently $\mathcal D$ has eight
invertibles and four simples of dimension $d$.

The local category $\mathcal K$ supplies all eight invertibles, while
$Z(\mathcal D)\simeq\mathcal B\boxtimes\mathcal K^{\mathrm{rev}}$
has sixteen. The kernel of the central forgetful map, and hence the
character group of the universal grading group $U$, has order two.
Equal dimensions of grading components, by comparing coefficients in
$a+cd^2=(16+32d)/|U|$, imply $c=4/|U|$. Thus $|U|$ divides four;
such a group is abelian, and therefore $U=C_2$.
The adjoint component has four invertibles and two noninvertibles.
Its two noninvertibles must map respectively to $X,Y$ under
de-equivariantization: two mapping to $X$ would contradict the occurrence
of $Y$ in $X^2$, and conversely. Their products force four distinct
invertible images. Hence $\mathcal D_{\rm ad}\to\mathcal C$ is a tensor
equivalence. The central boson $\ell$ lies outside this component and
gives the asserted tensor-product decomposition. This canonical
construction therefore supplies no exchange action producing a
near-group category of type $(8,16)$.
